% sec:dg-covariant — Covariant Derivatives and Christoffel Symbols
\section{Covariant Derivatives and Christoffel Symbols}
\label{sec:dg-covariant}

Differentiating a scalar field on a manifold is straightforward --- the partial
derivative $\pd{\mu} f$ transforms as a covector under coordinate changes, as
we saw in \cref{sec:dg-manifolds}.  But the ray tracer doesn't only work with
scalars.  It tracks a photon's four-momentum $p_\mu$ along a trajectory, and at
every step it must compute how that covector changes as the photon moves through
curved spacetime.  A naive partial derivative $\pd{\mu} V^\nu$ of a vector
field's components does not transform as a tensor: the coordinate basis vectors
themselves vary from point to point, and their variation contaminates the
derivative with non-tensorial terms.  The covariant derivative fixes this by
accounting for the twisting of the coordinate basis.

\paragraph{The failure of partial derivatives.}

To see the problem concretely, consider a vector field $V = V^\mu\,\pd{\mu}$ on
a manifold.  Under a coordinate change $x^\mu \to x'^\mu$, the components
transform as $V'^\mu = (\partial x'^\mu / \partial x^\nu)\, V^\nu$.  Taking
the partial derivative of both sides with respect to a new coordinate
$x'^\alpha$ yields
%
\begin{equation}
  \frac{\partial V'^\mu}{\partial x'^\alpha}
  = \frac{\partial x'^\mu}{\partial x^\nu}\,
    \frac{\partial V^\nu}{\partial x'^\alpha}
  + \frac{\partial^2 x'^\mu}{\partial x^\nu\,\partial x^\sigma}\,
    \frac{\partial x^\sigma}{\partial x'^\alpha}\, V^\nu \,.
  \label{eq:partial-nontensor}
\end{equation}
%
The first term is what a tensorial transformation would produce; the second term
--- involving the second derivative of the coordinate transformation --- is the
spoiler.  It vanishes only when the coordinate transformation is linear (as for
Lorentz transformations in flat spacetime), which is why partial derivatives
were perfectly adequate throughout \cref{ch:special-relativity}.  On a curved
manifold, no global linear coordinate system exists, and the second term is
generically non-zero.
\physnote{The non-tensorial term $\partial^2 x'^\mu / \partial x^\nu\,\partial
x^\sigma$ encodes how the coordinate grid itself curves.  The covariant
derivative cancels this term by subtracting out the basis variation.}

\paragraph{The covariant derivative.}

A \emph{covariant derivative} (or \emph{connection}) $\cd{\mu}$ is a map
that takes a tensor field and produces a new tensor field of one higher
covariant rank, while satisfying linearity, the Leibniz rule, and agreement
with partial derivatives on scalar fields ($\cd{\mu} f = \pd{\mu} f$).  For a
contravariant vector field, the covariant derivative takes the form
%
\begin{equation}
  \cd{\mu} V^\nu = \pd{\mu} V^\nu
    + \christoffel{\nu}{\mu}{\alpha}\, V^\alpha \,.
  \label{eq:covariant-deriv-vector}
\end{equation}
%
The additional term $\christoffel{\nu}{\mu}{\alpha}\, V^\alpha$ compensates for
the variation of the coordinate basis: it is precisely the non-tensorial
contribution that makes the combination transform as a $(1,1)$ tensor.  The
coefficients $\christoffel{\nu}{\mu}{\alpha}$ are called \emph{connection
coefficients} or, when derived from the metric, \emph{Christoffel symbols}.

For a covector field $\omega_\nu$, the corresponding formula has a minus sign
and the contraction is on a different index:
%
\begin{equation}
  \cd{\mu} \omega_\nu = \pd{\mu} \omega_\nu
    - \christoffel{\alpha}{\mu}{\nu}\, \omega_\alpha \,.
  \label{eq:covariant-deriv-covector}
\end{equation}
%
The sign flip and index placement follow from requiring the Leibniz rule to hold
for the contraction $\omega_\nu V^\nu$: differentiating a scalar must yield a
partial derivative, so the connection terms from the vector and covector must
cancel.
\warnnote{The plus sign in \cref{eq:covariant-deriv-vector} and the minus sign
in \cref{eq:covariant-deriv-covector} are a frequent source of sign errors.
The mnemonic: upper indices get a plus, lower indices get a minus.}

The extension to arbitrary tensors follows the same pattern.  For a $(1,1)$
tensor $T^\mu{}_\nu$, the covariant derivative is
%
\begin{equation}
  \cd{\alpha} T^\mu{}_\nu = \pd{\alpha} T^\mu{}_\nu
    + \christoffel{\mu}{\alpha}{\beta}\, T^\beta{}_\nu
    - \christoffel{\beta}{\alpha}{\nu}\, T^\mu{}_\beta \,.
  \label{eq:covariant-deriv-mixed}
\end{equation}
%
Each upper index contributes a $+\Gamma$ term; each lower index contributes a
$-\Gamma$ term.  This pattern generalises to tensors of any rank.

\paragraph{The Levi-Civita connection.}

The covariant derivative is not unique: any choice of connection coefficients
satisfying the algebraic requirements defines a valid covariant derivative.
General relativity singles out a unique connection by imposing two additional
conditions:
%
\begin{enumerate}
  \item \emph{Metric compatibility}: $\cd{\alpha}\, \metric = 0$.  Parallel
    transport preserves inner products --- lengths and angles are conserved
    along any curve.
  \item \emph{Torsion-free}: $\christoffel{\lambda}{\mu}{\nu} =
    \christoffel{\lambda}{\nu}{\mu}$.  The connection is symmetric in its lower
    two indices.
\end{enumerate}
%
The unique connection satisfying both conditions is the \emph{Levi-Civita
connection}, and its coefficients are the \emph{Christoffel symbols of the
second kind}.

\paragraph{Deriving the Christoffel symbols.}

The two conditions above determine the Christoffel symbols entirely in terms of
the metric.  The strategy is to write metric compatibility
$\cd{\alpha}\, g_{\mu\nu} = 0$ in component form and then exploit the
torsion-free symmetry to isolate $\Gamma$.  Expanding $\cd{\alpha}\,g_{\mu\nu}$
using \cref{eq:covariant-deriv-covector} on each lower index gives
%
\begin{equation}
  \pd{\alpha}\, g_{\mu\nu}
    = \christoffel{\lambda}{\alpha}{\mu}\, g_{\lambda\nu}
    + \christoffel{\lambda}{\alpha}{\nu}\, g_{\mu\lambda} \,.
  \label{eq:metric-compat-expanded}
\end{equation}
%
This is one equation relating partial derivatives of the metric to Christoffel
symbols.  To solve for $\Gamma$, write the same identity for two cyclic
permutations of the free indices:
%
\begin{align}
  \pd{\mu}\, g_{\nu\alpha}
    &= \christoffel{\lambda}{\mu}{\nu}\, g_{\lambda\alpha}
    + \christoffel{\lambda}{\mu}{\alpha}\, g_{\nu\lambda} \,,
  \label{eq:metric-compat-perm1} \\
  \pd{\nu}\, g_{\alpha\mu}
    &= \christoffel{\lambda}{\nu}{\alpha}\, g_{\lambda\mu}
    + \christoffel{\lambda}{\nu}{\mu}\, g_{\alpha\lambda} \,.
  \label{eq:metric-compat-perm2}
\end{align}
%
Adding these two equations and subtracting \cref{eq:metric-compat-expanded},
then using the torsion-free condition
$\christoffel{\lambda}{\mu}{\nu} = \christoffel{\lambda}{\nu}{\mu}$ to cancel
four of the six $\Gamma$ terms, leaves
$2\,\christoffel{\lambda}{\mu}{\nu}\, g_{\lambda\alpha}$ on the right.
Contracting with $g^{\alpha\sigma}$ and relabelling gives
%
\begin{equation}
  \christoffel{\lambda}{\mu}{\nu} = \frac{1}{2}\,g^{\lambda\sigma}\!
    \left(
      \pd{\mu}\, g_{\nu\sigma}
    + \pd{\nu}\, g_{\mu\sigma}
    - \pd{\sigma}\, g_{\mu\nu}
    \right) .
  \label{eq:christoffel-formula}
\end{equation}
%
This is the Christoffel symbol in terms of first derivatives of the metric and
the inverse metric.  Each index $\lambda,\mu,\nu$ runs from 0 to 3, so in four
dimensions there are $4 \times 10 = 40$ independent Christoffel symbols
(the factor of 10 comes from the symmetry in the lower pair).
\histnote{The notation $\christoffel{\lambda}{\mu}{\nu}$ for these coefficients
is due to Elwin Bruno Christoffel, who introduced them in 1869 --- decades
before Einstein's general relativity gave them a physical role.}

\paragraph{A worked example: the two-sphere.}

Returning to the two-sphere $S^2$ from \cref{sec:dg-metric}, with
the metric of \cref{eq:sphere-metric}, we compute the Christoffel symbols
from \cref{eq:christoffel-formula}.  Since the
metric is diagonal and $g_{\theta\theta} = R^2$ is constant, the only
components that survive are those involving $\pd{\theta}\, g_{\varphi\varphi}$.
As a concrete illustration, take $\christoffel{\theta}{\varphi}{\varphi}$:
%
\begin{align}
  \christoffel{\theta}{\varphi}{\varphi}
    &= \tfrac{1}{2}\, g^{\theta\theta}
      \bigl(\pd{\varphi}\,g_{\varphi\theta}
      + \pd{\varphi}\,g_{\varphi\theta}
      - \pd{\theta}\,g_{\varphi\varphi}\bigr)
  \notag \\
    &= \tfrac{1}{2} \cdot \frac{1}{R^2}
      \bigl(-2R^2 \sin\theta\cos\theta\bigr)
    = -\sin\theta\cos\theta \,.
  \label{eq:sphere-christoffel-detail}
\end{align}
%
The factor $R^2$ cancels between the metric derivative and the inverse metric.
Computing the remaining components similarly, the only non-vanishing Christoffel
symbols are
%
\begin{equation}
  \christoffel{\theta}{\varphi}{\varphi} = -\sin\theta\cos\theta \,,\qquad
  \christoffel{\varphi}{\theta}{\varphi}
    = \christoffel{\varphi}{\varphi}{\theta} = \cot\theta \,.
  \label{eq:sphere-christoffel}
\end{equation}
%
The Christoffel symbols of a sphere are independent of its radius --- they
encode the directions in which the basis vectors rotate, not the overall scale
of the geometry.

At the equator ($\theta = \pi/2$), all Christoffel symbols vanish: the
coordinate basis is parallel-transported along any direction through a point on
the equator, consistent with the equator being a geodesic of the sphere.
\physnote{The connection between Christoffel symbols and geodesics is made
precise by the geodesic equation in \cref{sec:dg-geodesic}.}

\paragraph{Metric compatibility in action.}

The condition $\cd{\alpha}\,\metric = 0$ has a direct computational
consequence: index raising and lowering commute with the covariant derivative.
That is, for any vector field $V^\mu$,
%
\begin{equation}
  \cd{\alpha}(g_{\mu\nu}\, V^\nu) = g_{\mu\nu}\, \cd{\alpha} V^\nu \,.
  \label{eq:metric-commutes-cd}
\end{equation}
%
The metric passes through the derivative like a constant.  This property is used
implicitly throughout the geodesic and ray-tracing chapters whenever we switch
between the momentum $p_\mu$ and the velocity
$\dot{x}^\mu = g^{\mu\alpha}\, p_\alpha$ --- the two formulations are
interchangeable precisely because the Levi-Civita connection preserves the
metric.
\xrefnote{The geodesic equation in \cref{sec:dg-geodesic} exploits metric
compatibility to move freely between Lagrangian ($\dot{x}^\mu$) and Hamiltonian
($p_\mu$) formulations.}

\paragraph{From explicit formulae to automatic differentiation.}

The Christoffel symbol formula (\cref{eq:christoffel-formula}) involves first
partial derivatives of the metric components --- 40 derivatives for a general
four-dimensional metric.  For a known analytic metric like Schwarzschild, one
can compute these by hand (or, more realistically, with a computer algebra
system).  For the Kerr metric, each of the 40 Christoffel symbols is a rational
function of $r$ and $\theta$ spanning several lines --- hand-differentiating and
transcribing them is error-prone and must be repeated for every new metric.  The
event-horizon codebase takes a different approach: it never computes Christoffel
symbols from hand-coded derivative expressions.
\coderef{Math.Tensor}{Christoffel4}
\implnote{The type \texttt{Christoffel4} stores four \texttt{Sym4 'Covariant}
values, one per upper index, exploiting the symmetry
$\christoffel{\lambda}{\mu}{\nu} = \christoffel{\lambda}{\nu}{\mu}$.  Each
\texttt{Sym4} stores 10 independent components.}

Instead, the codebase uses \emph{automatic differentiation}: the metric
function is written polymorphically (accepting any \texttt{Floating} type), and
the same code that evaluates $\metric$ at a point can compute exact
derivatives $\pd{\sigma}\,g_{\mu\nu}$ by passing an AD type instead of
\texttt{Double}.  The Christoffel formula (\cref{eq:christoffel-formula})
is then evaluated numerically at each integration step, using these exact
derivatives.  This strategy eliminates hand-coded derivative expressions,
reduces the surface area for implementation errors, and generalises
immediately to any metric the user supplies.
\xrefnote{The automatic differentiation strategy is developed in detail in
\cref{ch:autodiff}.}

The machinery is now in place to write down the equation that governs both
massive particles and photons in curved spacetime: the geodesic equation, which
states that freely falling worldlines have zero covariant acceleration.
